\chapter{联络、挠率与曲率}
\label{chap:ghs-manifold-connection}
\label{chap:ghs-geometry-foliation}

Lie 导数借助一个给定的流比较张量，但它不能回答另一类问题：如果没有指定流，怎样把一个切向量沿任意曲线搬到邻点，并定义“平行”与“加速度”？这需要额外的几何结构——联络。联络不是由光滑流形自动给出的；它规定不同切空间之间如何做无穷小比较。挠率与曲率随后不是独立添加的公式，而是联络的两个张量性缺陷：挠率比较协变微分的反对称部分与 Lie 括号，曲率比较沿两个方向连续平行移动的不交换性。

\section{向量丛联络与 仿射联络}
先给出比切丛更一般的母结构。设 $\pi:E\to M$ 为实或复向量丛。一个线性联络可以定义为
\begin{equation*}\nabla:\Gamma(E)\longrightarrow\Omega^1(M;E),
\qquad
\nabla(fs)=\dd f\otimes s+f\nabla s.
\end{equation*}
对向量场 $X$ 取收缩便得到 $\nabla_Xs:=(\nabla s)(X)$。在局部标架 $\{e_a\}$ 下，$\nabla e_b=\omega^a{}_b\otimes e_a$，因此任意截面 $s=s^ae_a$ 满足
\[
\nabla s=(\dd s^a+\omega^a{}_b s^b)\otimes e_a.
\]
联络的定义只依赖向量丛结构；切丛联络、法丛联络以及后面出现的主丛伴随丛联络都属于同一框架。

向量丛联络还可以从 主丛 的水平分布统一产生。设 $P\to M$ 是 主 $G$-丛，$\mathfrak g$ 是 $G$ 的 Lie 代数。对 $\xi\in\mathfrak g$，右作用生成 基本 竖直 矢量场
\[
\xi_P(p)
:=\left.\frac{\dd}{\dd t}\right|_0p\exp(t\xi).
\]
垂直子空间为
\[
V_pP=\ker\pi_{*p}
=\{\xi_P(p):\xi\in\mathfrak g\}.
\]
主 联络 可以等价地指定一个 $G$-equivariant 水平补空间
\begin{equation*}T_pP=H_pP\oplus V_pP,
\qquad
(R_g)_*H_pP=H_{pg}P,\end{equation*}
或者指定一个 $\mathfrak g$-值一形式 $\boldsymbol\omega\in\Omega^1(P;\mathfrak g)$，满足
\begin{equation}
\boldsymbol\omega(\xi_P)=\xi,
\qquad
R_g^*\boldsymbol\omega
=\operatorname{Ad}_{g^{-1}}\boldsymbol\omega.
\label{eq:principal-connection-form}
\end{equation}
此时
\[
H_pP=\ker\boldsymbol\omega_p.
\]
因此联络的几何本体是“怎样在总空间中选择水平提升”，局部矩阵值一形式只是选择截面后的表示。

若 $s_\alpha:U_\alpha\to P$ 是局部截面，定义局部 规范势
\[
A_\alpha=s_\alpha^*\boldsymbol\omega.
\]
在 $s_\beta=s_\alpha g_{\alpha\beta}$ 下，利用\eqnrefs{eq:principal-connection-form} 可得
\begin{equation}
A_\beta
=g_{\alpha\beta}^{-1}A_\alpha g_{\alpha\beta}
+g_{\alpha\beta}^{-1}\dd g_{\alpha\beta}.
\label{eq:principal-local-connection-transform}
\end{equation}
这就是\eqnrefs{eq:math-dg-connection-gauge-transform} 的 主-丛 来源。

给定表示 $\rho:G\to GL(W)$ 与 伴随丛 $E=P\times_\rho W$，主 联络 唯一诱导 $E$ 上的向量丛联络。局部上若截面写成 $W$-值函数 $\psi_\alpha$，则
\begin{equation*}D\psi_\alpha
=\dd\psi_\alpha+\rho_*(A_\alpha)\psi_\alpha.\end{equation*}
换 规范 后，$D\psi$ 按与 $\psi$ 相同的表示齐次变换。于是“非齐次变换的 联络 + 齐次变换的 物质 场”并不是物理中特有的记号，而是 伴随丛 的一般微分几何机制。

取 $E=TM$ 时，流形 $M$ 上的 仿射 联络 等价地写成双线性映射
\[
\nabla:\Gamma(TM)\times\Gamma(TM)\to\Gamma(TM),
\qquad
(X,Y)\mapsto\nabla_XY,
\]
满足对 $f\in C^\infty(M)$
\begin{align*}
\nabla_{fX}Y&=f\nabla_XY,\\
\nabla_X(fY)&=X(f)Y+f\nabla_XY.
\end{align*}
第一条说明它对“求导方向”是张量性的；第二条说明它对被求导向量满足 Leibniz 规则。与 Lie 导数不同，$\nabla_XY$ 对交换 $X,Y$ 没有预设关系。

取局部标架 $\{e_a\}$ 与余标架 $\{\theta^a\}$。定义联络系数
\begin{equation*}\nabla_{e_b}e_c=\Gamma^a{}_{bc}e_a,
\end{equation*}
并把同一信息组织成联络一形式
\begin{equation*}\nabla e_c=\omega^a{}_c\otimes e_a,
\qquad
\omega^a{}_c=\Gamma^a{}_{bc}\theta^b.
\end{equation*}
于是对 $Y=Y^ce_c$，
\[
\nabla_XY
=X(Y^a)e_a+Y^c\omega^a{}_c(X)e_a.
\]
这显示了联络一形式的角色：它补偿标架自身随位置的变化，使“分量求导 + 基底求导”重新组合成几何向量。

联络自然延伸到任意张量。对一形式 $\alpha$，定义
\[
(\nabla_X\alpha)(Y)
:=X[\alpha(Y)]-\alpha(\nabla_XY),
\]
再要求与张量积和缩并相容，就唯一确定一般张量的协变导数。对含标架指标的微分形式还需要把普通外微分升级为协变外微分。若 $\eta=\eta^ae_a$ 是向量值 $p$-形式，定义
\begin{equation}\label{eq:math-dg-exterior-covariant-derivative}
(D\eta)^a
:=\dd\eta^a+\omega^a{}_b\wedge\eta^b.
\end{equation}
若指标为协变指标，相应联络项带负号；多个指标时对每个指标分别作用。$D$ 是普通 $\dd$ 在“带联络的向量丛值形式”上的延伸，后面的 Bianchi 恒等式都用这一定义。

联络系数不是张量。若换标架 $e'=e\Lambda$，其中 $\Lambda:U\to GL(n,\RR)$，则
\begin{equation}\label{eq:math-dg-connection-gauge-transform}
{\omega'=\Lambda^{-1}\omega\Lambda+\Lambda^{-1}\dd\Lambda.}
\end{equation}
非齐次项 $\Lambda^{-1}\dd\Lambda$ 正是“联络不是张量”的精确含义。

向量丛值形式上的曲率最简洁地由 $D^2$ 定义。对 $E$-值 $p$-形式 $\eta=\eta^ae_a$，令
\[
\Omega:=\dd\omega+\omega\wedge\omega.
\]
则有
\begin{equation}\label{eq:math-dg-covariant-derivative-square}
{D^2\eta=\Omega\wedge\eta.}
\end{equation}
这条公式先于切丛的指标表达成立；对 $E=TM$ 才把 $\Omega$ 解释为 Riemann 曲率二形式。

联络本身不是向量空间中的元素：两个联络可以相减，但单独一个联络不能在不选参考点的情况下被称为“零”。设 $\nabla$ 与 $\widetilde\nabla$ 是同一向量丛 $E$ 上的两个联络，定义
\[
\mathcal A_Xs
:=(\widetilde\nabla_X-\nabla_X)s.
\]
Leibniz 项精确相消，因此 $\mathcal A_Xs$ 对 $X$ 与 $s$ 都是 $C^\infty(M)$-线性的。故
\begin{equation*}{
\widetilde\nabla-\nabla
\in\Omega^1(M;\End E).
}\end{equation*}
反过来，任意 $A\in\Omega^1(M;\End E)$ 都定义新联络
\[
(\nabla+A)_Xs:=\nabla_Xs+A(X)s.
\]
所以所有联络组成一个以 $\Omega^1(M;\End E)$ 为模型向量空间的 仿射 空间。

在局部标架中，若参考联络形式为 $\omega$，则新联络形式只是
\[
\widetilde\omega=\omega+A.
\]
但曲率不是线性变化。由
\[
F_\nabla=\dd\omega+\omega\wedge\omega
\]
直接展开得
\begin{equation}
{
F_{\nabla+A}
=F_\nabla+D_\nabla A+A\wedge A,
}
\label{eq:math-dg-curvature-finite-perturbation}
\end{equation}
其中
\[
D_\nabla A
:=\dd A+\omega\wedge A+A\wedge\omega
\]
是 $\End(E)$-值一形式上的协变外微分。若取一参数族 $\nabla_t=\nabla+tA$，则一阶变分为
\begin{equation*}\left.\frac{\dd}{\dd t}\right|_{t=0}F_{\nabla_t}
=D_\nabla A.\end{equation*}
Chern--Weil 超渡、Yang--Mills 线性化以及后面度量变分中的 $\partial_tR$ 都是这一一般公式的不同表现。

对切丛 仿射 联络，差张量可写成
\[
A(X,Y):=(\widetilde\nabla_X-\nabla_X)Y
\in\Gamma(TM).
\]
挠率的变化特别简单：
\begin{equation*}{
\widetilde T(X,Y)
=T(X,Y)+A(X,Y)-A(Y,X).
}\end{equation*}
所以给定一个联络后，只需调整差张量的反对称部分就能改变 挠率；其对称部分不影响 挠率。曲率则满足有限差分公式
\begin{align}
\widetilde R(X,Y)Z
={}&R(X,Y)Z
+(\nabla_XA)(Y,Z)-(\nabla_YA)(X,Z)\notag\\
&+A(X,A(Y,Z))-A(Y,A(X,Z))\notag\\
&+A(T(X,Y),Z),
\label{eq:math-dg-curvature-connection-change}
\end{align}
其中最后一项反映参考联络可能本身有 挠率；若 $T=0$ 就消失。

若 $E$ 带 纤维 度量 $h$，保持 $h$ 的联络差张量必须取值于 $h$-反自伴 自同态；在正交标架下就是 $\mathfrak{so}(r)$-值一形式。特别地，对 Riemann 切丛，度量-兼容 联络 构成由 $\Omega^1(M;\mathfrak{so}(n))$ 平移得到的 仿射 子空间。下一章再加入 挠率-自由 条件后，这个 仿射 自由度 被完全固定，得到唯一 Levi--Civita 联络。

\section{挠率、曲率与 Cartan 结构方程}
联络给定后，先比较 $\nabla_XY-\nabla_YX$ 与光滑结构已经给出的 $[X,Y]$。它们的差定义挠率
\begin{equation*}T(X,Y)
:=\nabla_XY-\nabla_YX-[X,Y].
\end{equation*}
逐项检查可知对 $X,Y$ 都是 $C^\infty(M)$-线性的，因此 $T$ 是真正的 $(1,2)$ 张量。写
\[
T(X,Y)=T^a(X,Y)e_a
\]
定义挠率二形式 $T^a$。

曲率则比较沿 $X,Y$ 两个方向连续协变微分的交换子，并扣除方向场本身的不交换：
\begin{equation*}R(X,Y)Z
:=\nabla_X\nabla_YZ-\nabla_Y\nabla_XZ-\nabla_{[X,Y]}Z.
\end{equation*}
它对 $X,Y,Z$ 都是函数线性的，因此是 $(1,3)$ 张量。定义曲率二形式 $\Omega^a{}_b$ 为
\[
R(X,Y)e_b=\Omega^a{}_b(X,Y)e_a.
\]
若
\[
R(e_c,e_d)e_b=R^a{}_{bcd}e_a,
\]
则
\[
\Omega^a{}_b=\frac12R^a{}_{bcd}\theta^c\wedge\theta^d.
\]

Cartan 结构方程由此写成外微分形式：挠率 $T^a=\dd\theta^a+\omega^a{}_b\wedge\theta^b$（\eqnrefs{eq:math-dg-synthesis-cartan-first}），曲率 $\Omega^a{}_b=\dd\omega^a{}_b+\omega^a{}_c\wedge\omega^c{}_b$（\eqnrefs{eq:math-dg-synthesis-cartan-second}）。

在一般标架中，把 $\omega^a{}_c=\Gamma^a{}_{bc}\theta^b$ 与
$[e_b,e_c]=C^a{}_{bc}e_a$ 代入第一结构方程，挠率分量为
\begin{equation*}T^a{}_{bc}
=\Gamma^a{}_{bc}-\Gamma^a{}_{cb}-C^a{}_{bc}.
\end{equation*}
因此“联络系数下标对称”只在坐标基 $C^a{}_{bc}=0$ 中才与无挠等价。非坐标基中若忘掉结构函数，会直接得到错误结论。

在 主丛 语言中，曲率形式为
\begin{equation*}\mathbf F
=\dd\boldsymbol\omega
+\frac12[\boldsymbol\omega\wedge\boldsymbol\omega],\end{equation*}
拉回局部截面后成为
\begin{equation*}F=\dd A+A\wedge A.\end{equation*}
换 规范 时没有非齐次项：
\begin{equation*}F'=g^{-1}Fg.\end{equation*}
因此曲率是 tensorial object；它真正测量联络本身，而不是局部 标架 的选择。

若 $F=0$，称联络 平坦。单连通坐标域上，平坦性正是局部存在 规范 $g$ 使
\begin{equation*}A=g^{-1}\dd g\end{equation*}
的可积条件；再作 规范变换 $g^{-1}$ 就可令局部 联络形式 消失。于是
\[
F=0
\quad\Longleftrightarrow\quad
A\text{ 局部可 规范 到 }0
\]
在单连通局部成立。但全局平坦联络仍可能具有非平凡 和乐，因为 noncontractible 回路 能保存全局拓扑信息。

曲率二形式的换标架律则是齐次的：
\begin{equation*}\Omega'=\Lambda^{-1}\Omega\Lambda.
\end{equation*}
这表明曲率本身是几何张量，虽然表示它的联络一形式不是。

\section{Bianchi 恒等式}
Cartan 方程并不是任意二形式方程。因为 $T$ 与 $\Omega$ 都来自同一个联络，它们自动满足两条微分一致性条件。

由此得到结构方程的 Bianchi 恒等式：$DT^a=\Omega^a{}_b\wedge\theta^b$ 与 $D\Omega^a{}_b=0$（\eqnrefs{eq:math-dg-bianchi-first}、\eqnrefs{eq:math-dg-bianchi-second}、\eqnrefs{eq:math-dg-synthesis-bianchi}）。

\par
\begin{workedexample}
从 Cartan 第一、第二结构方程出发，不使用分量公式，推导两条 Bianchi 恒等式。

\textbf{解。} 记 \(D\) 为由联络诱导的协变外微分。对 \(T^a=D\theta^a\) 再作用 \(D\)，利用 \(D^2\theta^a=\Omega^a{}_b\wedge\theta^b\)，得
\[
DT^a=\Omega^a{}_b\wedge\theta^b.
\]
对曲率定义 \(\Omega=\dd\omega+\omega\wedge\omega\) 作用协变外微分，展开后由 \(\dd^2=0\) 和结合律消去项，得 \(D\Omega^a{}_b=0\)。这正是\eqnrefs{eq:math-dg-bianchi-first} 与\eqnrefs{eq:math-dg-bianchi-second}；分量展开见本章末推导框\eqnrefs{eq:math-dg-synthesis-bianchi}。
\end{workedexample}
\par


第二 Bianchi 恒等式的一个深层后果是：对任意 Ad-不变 多项式 $P$ on $\mathfrak g$，把曲率代入所得微分形式
\begin{equation*}P(F)\in\Omega^{2k}(M)\end{equation*}
自动闭合。最常见的矩阵表示例子是
\[
\operatorname{tr}(F^k).
\]
由于 迹 消去 对易子，
\[
\dd\operatorname{tr}(F^k)
=\operatorname{tr}(D F^k)
=k\operatorname{tr}((DF)F^{k-1})=0.
\]
因此它定义 de Rham 上同调 类。

更重要的是，这个 上同调类 不依赖具体 联络。令 $A_t$ 是一族联络，$F_t$ 对应曲率。则
\begin{equation}
\frac{\dd}{\dd t}\operatorname{tr}(F_t^k)
=k\,\dd\operatorname{tr}
\left(\dot A_t\wedge F_t^{k-1}\right).
\label{eq:chern-weil-transgression}
\end{equation}
右端是 恰当形式，所以 $[\operatorname{tr}(F_t^k)]$ 与 $t$ 无关。曲率虽然依赖联络，某些由曲率构造的上同调类却只依赖 丛 拓扑；这就是 Chern--Weil 理论的核心机制。

对复向量丛，适当归一化的
\[
\det\left(I+\frac{\ii}{2\pi}F\right)
=1+c_1(E)+c_2(E)+\cdots
\]
给出 Chern 类 的 de Rham 代表元；对实定向丛则由 曲率 的 不变 多项式 构造 Pontryagin 类 与 Euler 类。给定向量丛后可以选择联络，由联络计算曲率；将曲率代入适当的不变多项式得到闭形式。联络变化只改变这些形式中的恰当项，所以它们代表丛本身的拓扑上同调类。
这也是为什么曲率既是局部平行移动缺陷，又能够携带全局拓扑信息。

第一 Bianchi 恒等式把挠率与曲率联系起来；第二 Bianchi 恒等式表达曲率自身的协变闭合。它们都不是运动方程，而是由联络定义强制产生的恒等式。


联络已由一般几何定义。为了将它用于场方程，需要把坐标基底的变化、体积因子和真正的曲率分别算出。





平直空间的径向量子化从一族同心球面开始。但球面不是量子化的本质：重要的是，能否用一族互不相交的超曲面逐次读取场的数据。对于一般空间，既不能先写 $r=|x|$，也不能默认这些超曲面具有相同的几何。本章从流形本身推出这些问题的正确表述。

\subsection{联络与曲率}
\label{sec:ghs-connection-curvature}

在平直空间的直角坐标中，我们往往不区分向量的普通导数和协变导数。可是只要在同一个平面上改用极坐标，这种区分就不得不出现。几何上的联络首先解决的是一个朴素问题：\emph{不同位置的向量属于不同的切空间，怎样比较它们的变化？}

考虑二维欧氏平面，先用直角坐标 $(X,Y)$ 写度规 $\dd s^2=\dd X^2+\dd Y^2$。引入 $X=r\cos\theta$、$Y=r\sin\theta$，微分为
\[
 \dd X=\cos\theta\,\dd r-r\sin\theta\,\dd\theta,\qquad
 \dd Y=\sin\theta\,\dd r+r\cos\theta\,\dd\theta.
\]
把两式平方相加，混合项相消，得到 $\dd s^2=\dd r^2+r^2\dd\theta^2$。这时 $g_{rr}=1$、$g_{\theta\theta}=r^2$，体积元为 $r\dd r\dd\theta$。$r$ 的因子是坐标面积伸缩，不是曲率：平面仍然是平面。

一般坐标中的切向基底是 $\partial_a$。向量 $V=V^a\partial_a$ 的普通分量导数 $\partial_bV^a$ 不能反映基底 $\partial_a$ 本身的变化。定义联络系数使
\begin{equation}\label{eq:ghs-bridge-christoffel-definition}
 \nabla_{\partial_b}\partial_c=\Gamma^a{}_{bc}\partial_a,
 \qquad\nabla_bV^a=\partial_bV^a+\Gamma^a{}_{bc}V^c.
\end{equation}
如果联络不引入新的扭率，并要求平行移动保持向量夹角，那么这两个条件分别写为 $\Gamma^a{}_{bc}=\Gamma^a{}_{cb}$、$\nabla_ag_{bc}=0$。后一条件展开为
\[
 \partial_ag_{bc}=\Gamma^d{}_{ab}g_{dc}+\Gamma^d{}_{ac}g_{bd}.
\]
将指标 $(a,b,c)$ 的三个循环形式作两次加、一次减，便得到
\begin{equation}\label{eq:ghs-bridge-Christoffel-metric}
 \Gamma^a{}_{bc}=\frac12g^{ad}
 (\partial_bg_{dc}+\partial_cg_{db}-\partial_dg_{bc}).
\end{equation}
这个公式不是从某个特殊坐标猜出来的，而是由\emph{度规相容性和零扭率}唯一决定。

对平面的极坐标代入 $\partial_rg_{\theta\theta}=2r$，立即得到
\[
 \Gamma^r{}_{\theta\theta}=-r,\qquad
 \Gamma^\theta{}_{r\theta}=\Gamma^\theta{}_{\theta r}=\frac1r,
\]
其余分量为零。由此可见，$\Gamma$ 的某个分量不为零不等于空间弯曲；它可能仅仅反映坐标基底随位置旋转。

矢量场 $V^a$ 的散度可以通过积分定理决定。要求 $\int_U\operatorname{div}V\,\dd\mu_g$ 等于边界的法向通量，会得到
\[
 \operatorname{div}_gV=\frac1{\sqrt{g}}\partial_a(\sqrt g\,V^a),
 \qquad g=\det(g_{ab}).
\]
标量函数 $f$ 的梯度满足 $(\operatorname{grad}f)^a=g^{ab}\partial_bf$。先取梯度、再取散度，得到
\begin{equation}\label{eq:ghs-bridge-laplace-beltrami}
 \Delta_gf=\frac1{\sqrt g}\partial_a\bigl(\sqrt g\,g^{ab}\partial_bf\bigr).
\end{equation}
在二维极坐标中，代入 $\sqrt g=r$、$g^{rr}=1$、$g^{\theta\theta}=r^{-2}$，可逐项算出
\[
 \Delta_gf=\frac1r\partial_r(r\partial_rf)
 +\frac1r\partial_\theta\left(\frac1r\partial_\theta f\right)
 =\partial_r^2f+\frac1r\partial_rf+\frac1{r^2}\partial_\theta^2f.
\]
从 \cref{eq:ghs-bridge-laplace-beltrami} 走到这里的每一个步骤只用链式法则和分部积分。后文的曲空间自由场，无论几何多么复杂，都采用同一个定义；不同之处是 $g_{ab}$ 不再能在整个区域变成欧氏度规。

当向量沿一个无穷小闭合回路平行移动时，比较先沿 $a$ 再沿 $b$ 与相反顺序的结果，会得到
\[
 [\nabla_a,\nabla_b]V^c=R^c{}_{dab}V^d,
\]
其中曲率可由联络计算：
\begin{equation}\label{eq:ghs-bridge-Riemann-coordinate}
 R^a{}_{bcd}=\partial_c\Gamma^a{}_{db}-\partial_d\Gamma^a{}_{cb}
 +\Gamma^a{}_{ce}\Gamma^e{}_{db}-\Gamma^a{}_{de}\Gamma^e{}_{cb}.
\end{equation}
在极坐标平面上，例如 $R^r{}_{\theta r\theta}$ 的导数项为
$\partial_r(-r)=-1$，而非零二次联络项为
$-\Gamma^r{}_{\theta\theta}\Gamma^\theta{}_{r\theta}=+1$，两者正好抵消。这个计算显示：真正的曲率要由联络的导数和乘积共同决定，而不能单看 $\Gamma$。

对半径 $R_0$ 的圆球面，度规为 $\dd s^2=R_0^2(\dd\vartheta^2+\sin^2\vartheta\,\dd\phi^2)$。代入 \cref{eq:ghs-bridge-Christoffel-metric}，得到 $\Gamma^\vartheta{}_{\phi\phi}=-\sin\vartheta\cos\vartheta$、$\Gamma^\phi{}_{\vartheta\phi}=\cot\vartheta$。再代入曲率式可得 $R^\vartheta{}_{\phi\vartheta\phi}=\sin^2\vartheta$，于是 $\mathcal R=2/R_0^2$。球面曲率不会像平面的极坐标联络那样完全抵消。这两个算例足以使联络、曲率和坐标效应之间的区别变得可计算，而不需要先背诵抽象微分几何定理。

设 $M$ 是连通的 $d$ 维光滑流形，$d\geq 2$，$g$ 是正定黎曼度规，$\nabla$ 为其 Levi--Civita 联络。在局部坐标 $x^\mu$ 中，体积测度是 $\dd\mu_g=\sqrt{\det g_{\mu\nu}}\,\dd^d x$。坐标的存在是局部事实；是否有整个空间通用的径向坐标是额外的全局问题。

\subsection{关联张量丛与高阶协变导数}

局部标架表示规定场的纤维类型；联络将这一表示变成跨点的微分运算。

前面的群表示分解描述的是全局等距变换。若要研究场本身的自旋，还必须回到每一点的切空间，弄清张量怎样在不同点之间进行比较。

令 $\operatorname{Fr}_{SO}(M)$ 为定向正交标架主丛，结构群 $SO(d)$。给定其有限维表示 $\rho:SO(d)\to GL(V)$，定义关联向量丛
\begin{equation}\label{eq:ghs-associated-vector-bundle}
 E_\rho=\operatorname{Fr}_{SO}(M)\times_\rho V.
\end{equation}
它的截面 $\Gamma(E_\rho)$ 是空间上的场。对于整数自旋 $S$ 的对称无迹场，可取
\[
 E_S=\operatorname{Sym}^S_0(T^*M),\qquad
 \mathcal J^{(S)}\in\Gamma(E_S).
\]
当 $d>2$ 时，纤维维数为 $\binom{d+S-1}{S}-\binom{d+S-3}{S-2}$，第二项在 $S<2$ 时取零；这是对称张量空间减去迹像的维数，《群论》的调和张量章节给出完整分解，且与位置无关。物理自由度是否等于这一纤维维数，仍取决于运动方程、约束和规范对称性。

自旋量场需要更严格的前提：存在自旋结构时，正交标架丛才能提升到 $\operatorname{Spin}(d)$ 主丛。并非每个定向流形都有自旋结构；其障碍由第二 Stiefel--Whitney 类 $w_2(M)$ 表征。因此“任意空间上的任意自旋场”不能在不区分整数自旋与半整数自旋的情况下成立。

Levi--Civita 联络自然延拓到所有张量丛。对向量场 $V^c$，我们规定曲率号约定
\begin{equation}\label{eq:ghs-curvature-commutator}
 [\nabla_a,\nabla_b]V^c=R^c{}_{dab}V^d.
\end{equation}
对协向量使用度规相容性与 Leibniz 规则可得
$[\nabla_a,\nabla_b]\omega_c=-R^d{}_{cab}\omega_d$。推广到秩 $S$ 协变张量时，右端是 $S$ 个曲率项之和；每个指标单独贡献一次曲率作用。

\begin{proposition}[标量二阶导数与三阶导数]
对任意标量 $f$，$[\nabla_a,\nabla_b]f=0$；但在约定 \cref{eq:ghs-curvature-commutator} 下
\begin{equation}\label{eq:ghs-scalar-third-derivative}
 (\nabla_a\nabla_b-\nabla_b\nabla_a)(\nabla_cf)
 =-R^d{}_{cab}\nabla_df.
\end{equation}
\end{proposition}
\begin{proof}
标量的一阶协变导数就是偏导，故前两次作用交换。其梯度却是协向量，于是第三次求导时按照协向量的曲率对易关系获得右端。
\end{proof}

\begin{figure}[htbp]
\centering
% Local TikZ geometry; styles inherited exclusively from the shared preamble.
\begin{tikzfig}[scale=1.05]
\coordinate (O) at (0,0);\coordinate (A) at (2.65,.35);\coordinate (B) at (3.4,1.85);\coordinate (C) at (.75,1.5);
\fill[gray!7] (O)--(A)--(B)--(C)--cycle;
\draw[thick,lecturearrow] (O)--(A) node[midway,below] {$u$};
\draw[thick,lecturearrow] (A)--(B) node[midway,right] {$v$};
\draw[thick,lecturearrow] (B)--(C);
\draw[thick,lecturearrow] (C)--(O);
\fill (O) circle (1.6pt);
\draw[blue!65!black,lecturearrow] (O)--(.13,.9) node[left] {$V$};
\draw[red!65!black,lecturearrow] (O)--(.4,.81) node[right] {$V+\delta V$};
\node[lecturelabel] at (1.6,-.8) {$\delta V\propto R(u,v)V\;\cdot\;\mathrm{Area}$};
\end{tikzfig}

\caption{小闭合回路上的平行移动示意。沿两个方向交换传输顺序后，向量的差在最低非平凡阶由曲率 $R(u,v)$ 乘以有向面积控制。}\label{fig:ghs-curvature-holonomy}
\end{figure}

\cref{fig:ghs-curvature-holonomy} 中画出的两条末态箭头不是额外的自由度，而是同一向量在不同微小路径下平行移动的比较；当曲率为零时，它们在最低阶重合。

这说明“把平直空间公式中的所有 $\partial$ 换成 $\nabla$”不保证新公式守恒。高自旋流经常包含三阶及更多阶导数，曲率对易项在散度计算中不可避免。在一般曲率背景上，须同时考虑曲率修正、背景方程及可能的几何障碍。


\section{平行移动与曲率的几何意义}
沿光滑曲线 $\gamma(t)$，向量场 $V(t)\in T_{\gamma(t)}M$ 若满足
\begin{equation*}\nabla_{\dot\gamma}V=0,
\end{equation*}
则称 $V$ 沿 $\gamma$ 平行。写 $V=V^ae_a$ 后，这是线性常微分方程
\[
\frac{\dd V^a}{\dd t}
+\omega^a{}_b(\dot\gamma)V^b=0.
\]
因此给定联络以后，任意曲线都诱导一个线性平行移动映射。

局部标架中，平行移动方程可写成矩阵 ODE
\begin{equation}
\dot V(t)+A_tV(t)=0,
\qquad
A_t:=\omega(\dot\gamma(t)).
\label{eq:parallel-matrix-ode}
\end{equation}
若不同时刻的 $A_t$ 彼此对易，则
\[
V(t)=\exp\left(-\int_0^tA_\tau\,\dd\tau\right)V(0).
\]
一般情况下必须保留时间次序，写成 path-ordered 指数
\begin{equation*}U_\gamma(t,0)
=\mathcal P\exp\left(-\int_0^tA_\tau\,\dd\tau\right).\end{equation*}
其 Dyson 级数 为
\[
I+
\sum_{k=1}^\infty(-1)^k
\int_{0<t_k<\cdots<t_1<t}
A_{t_1}\cdots A_{t_k}
\,\dd t_k\cdots\dd t_1.
\]
所以 path ordering 不是记号装饰，而是非交换 联络 矩阵 的精确解法。

若 $\gamma$ 是闭路，$U_\gamma$ 给出 和乐 变换。对面积为 $\epsilon^2$ 的微小平行四边形，其 和乐 有展开
\begin{equation*}U_{\partial\Sigma}
=I-
F(X,Y)\epsilon^2+O(\epsilon^3),\end{equation*}
符号取决于平行移动方程和闭路定向约定。核心内容是二阶系数恰为 曲率。后续 和乐 章节会把这个局部关系提升为 和乐 代数 与曲率生成元的全局定理。

曲率的局部意义可以从无穷小闭合回路看出。若先沿 $X$ 再沿 $Y$ 平行移动，并与相反次序比较，扣除底点回路本身由 $[X,Y]$ 造成的偏移以后，最低非平凡差异正由 $R(X,Y)$ 给出。换言之，$R=0$ 意味着局部平行移动对足够小的可缩闭路没有二阶 和乐；$R\neq0$ 则说明不同路径的比较无法彼此兼容。
联络的一般定义给出了求导规则。以下从该规则推导局部系数和曲率分量，并在活动标架中检查同一个几何对象怎样改写。

\section{联络系数从定义推出}
\label{sec:math-connection-coefficients}

对切丛 \(E=TM\)，在坐标基 \(\partial_i\) 上定义函数
\(\Gamma^k{}_{ij}\)：
\[
\nabla_{\partial_i}\partial_j=\Gamma^k{}_{ij}\partial_k.
\]
若 \(X=X^i\partial_i\)、\(Y=Y^j\partial_j\)，联络对第一个变量的
函数线性性及对第二个变量的乘积法则给
\begin{equation}
(\nabla_XY)^k
=X^i\bigl(\partial_iY^k+\Gamma^k{}_{ij}Y^j\bigr).
\label{eq:math-covariant-derivative-coordinates}
\end{equation}
第一个加数对坐标分量求导，第二个加数补偿坐标基随位置的变化。

为什么 \(\Gamma^k{}_{ij}\) 不是张量？令新坐标为
\(y^a=y^a(x)\)。前章已证明
\(\partial/\partial y^a=(\partial x^i/\partial y^a)\partial_i\)。
把它代入 \(\nabla_{\partial/\partial y^a}(\partial/\partial y^b)\)：
第二个基是 \((\partial x^j/\partial y^b)\partial_j\)，其系数随位置变化，
故乘积法则给出
\[
\nabla_{\frac{\partial}{\partial y^a}}
\left(\frac{\partial x^j}{\partial y^b}\partial_j\right)
=
\frac{\partial^2 x^j}{\partial y^a\partial y^b}\,\partial_j
+
\frac{\partial x^i}{\partial y^a}\frac{\partial x^j}{\partial y^b}
\Gamma^k{}_{ij}\,\partial_k .
\]
第一项正是坐标基本身随位置转动产生的二阶导数；再把输出基
\(\partial_j,\partial_k\) 都换成 \(\partial/\partial y^c\)
（两边乘 \(\partial y^c/\partial x^k\)），得到
\begin{equation}
\Gamma'^c{}_{ab}
=\frac{\partial y^c}{\partial x^k}
\left(
\frac{\partial x^i}{\partial y^a}
\frac{\partial x^j}{\partial y^b}\Gamma^k{}_{ij}
+\frac{\partial^2x^k}{\partial y^a\partial y^b}
\right).
\label{eq:math-connection-transform}
\end{equation}
括号最后一项含坐标变换的二阶导数，张量的变换律没有这种项。
把\cref{eq:math-connection-transform} 代回
\cref{eq:math-covariant-derivative-coordinates}，它恰消掉
\(\partial_iY^k\) 变换时出现的多余项，使 \(\nabla_XY\) 真正按向量变换。

平面是检验该式的简单反例。在直角坐标中取通常平直联络，
\(\Gamma^k{}_{ij}=0\)。改用极坐标 \((r,\theta)\)，其中
\(\vb{e}_r=(\cos\theta,\sin\theta)\)、
\(\vb{e}_\theta=(-\sin\theta,\cos\theta)\) 是 \(\mathbb R^2\) 里
单位长正交向量，而坐标基 \(\partial_\theta=r\vb{e}_\theta\) 长 \(r\) 倍。
直接对坐标基求导，得到
\[
\nabla_{\partial_\theta}\partial_\theta=-r\partial_r,\qquad
\nabla_{\partial_r}\partial_\theta
=\nabla_{\partial_\theta}\partial_r=r^{-1}\partial_\theta.
\]
故 \(\Gamma^r{}_{\theta\theta}=-r\)、
\(\Gamma^\theta{}_{r\theta}=\Gamma^\theta{}_{\theta r}=1/r\)，
即使平面本身没有弯曲，联络系数也可以非零。

\section{平行移动是可解的初值问题}

给定曲线 \(\gamma(t)\) 与沿曲线的向量 \(V(t)=V^k(t)\partial_k\)，
把它代入\cref{eq:math-covariant-derivative-coordinates}，得到
\[
\nabla_{\dot\gamma}V=0
\quad\Longleftrightarrow\quad
\dot V^k+\Gamma^k{}_{ij}(\gamma(t))
\dot\gamma^i V^j=0.
\]
右边是一组线性常微分方程。若系数沿曲线连续，给定
\(V(t_0)\in T_{\gamma(t_0)}M\) 后，在曲线存在的区间中局部有唯一解。
这样联络把起点的向量搬到邻近点，且无需先把两个切空间直接相减。

若联络还满足度量相容条件
\(X[g(Y,Z)]=g(\nabla_XY,Z)+g(Y,\nabla_XZ)\)，
则两个平行向量 \(V,W\) 沿同一曲线满足
\[
\frac{\dd}{\dd t}g(V,W)
=g(\nabla_{\dot\gamma}V,W)+g(V,\nabla_{\dot\gamma}W)=0.
\]
也就是说，长度和夹角在平行移动中保持；下一节将问什么条件使这样的
联络由度量唯一决定。

把前面的平面极坐标例子算到底，能避免一种常见误解。
沿半径为 \(r_0\) 的圆周取 \(\gamma(t)=(r_0,t)\)，
\(0\le t\le2\pi\)。平行方程成为
\[
\dot V^r-r_0V^\theta=0,
\qquad
\dot V^\theta+r_0^{-1}V^r=0.
\]
第二个坐标基向量 \(\partial_\theta\) 的长度是 \(r_0\)，
所以令 \(W^\theta=r_0V^\theta\) 才得到正交基中的物理分量。
此时 \(\dot V^r=W^\theta\)、\(\dot W^\theta=-V^r\)。
若起点分量为 \((A,B)\)，解为
\[
V^r(t)=A\cos t+B\sin t,
\qquad
W^\theta(t)=-A\sin t+B\cos t.
\]
走满一圈后分量回到 \((A,B)\)。沿途联络系数非零，
但闭路平行移动的净结果仍是恒等；这与平面曲率为零一致。



\par
\begin{workedexample}
在同一平面取不是闭合圆周的弧线 \(0\le t\le\pi/2\)，
将上式的正交基分量与直角坐标分量比较。
解释为什么正交基分量旋转了 \(\pi/2\)，平行移动的向量
在固定直角坐标中却没有旋转。
\textbf{解.} 对圆弧 \(\gamma(t)=(r_0,t)\)，已有平行解的正交分量为
\[
V^r=A\cos t+B\sin t,\qquad W^\theta=-A\sin t+B\cos t.
\]
在 \(t=\pi/2\) 它由 \((A,B)\) 变为 \((B,-A)\)，即在转动的极坐标正交标架中旋转 \(-\pi/2\)。但 \(e_r=(\cos t,\sin t)\)、\(e_\theta=(-\sin t,\cos t)\)，故固定笛卡尔分量
\[
V^x=V^r\cos t-W^\theta\sin t=A,\qquad
V^y=V^r\sin t+W^\theta\cos t=B
\]
保持不变。变化来自标架转动，不是向量在平面中转动。
\end{workedexample}
\par


\section{曲率分量从两次协变求导得到}
\label{sec:math-curvature-components}

在坐标基中，\([\partial_i,\partial_j]=0\)。令
\(R(\partial_i,\partial_j)\partial_k
=R^\ell{}_{kij}\partial_\ell\)。
对 \(\nabla_{\partial_j}\partial_k
=\Gamma^\ell{}_{jk}\partial_\ell\) 再沿 \(\partial_i\) 求导，
分别对系数和基使用联络的乘积法则，便得
\begin{equation}
R^\ell{}_{kij}
=\partial_i\Gamma^\ell{}_{jk}
-\partial_j\Gamma^\ell{}_{ik}
+\Gamma^m{}_{jk}\Gamma^\ell{}_{im}
-\Gamma^m{}_{ik}\Gamma^\ell{}_{jm}.
\label{eq:math-curvature-components}
\end{equation}
每项都有来处：前两项是联络系数随位置的变化，后两项是
第一次求导后产生的基向量还要再次求导。

平面极坐标能检验“联络系数非零不代表曲率非零”。
上一节得到 \(\Gamma^r{}_{\theta\theta}=-r\)、
\(\Gamma^\theta{}_{r\theta}
=\Gamma^\theta{}_{\theta r}=1/r\)。
由\cref{eq:math-curvature-components}，
\[
R^r{}_{\theta r\theta}
=\partial_r(-r)
-\Gamma^\theta{}_{r\theta}\Gamma^r{}_{\theta\theta}
=-1-\frac1r(-r)=0.
\]
其余独立分量也为零；平面极坐标中的非零 \(\Gamma\)
仅反映基向量随位置改变。


\par
\begin{workedexample}
把单位球面的度量整体乘 \(a^2\)，其中 \(a>0\)。
从新的正交余标架求 \(\omega^1{}_2\) 与 \(\Omega^1{}_2\)，
证明半径为 \(a\) 的球面 Gauss 曲率为 \(1/a^2\)。
\textbf{解.} 取正交余标架 \(\theta^1=a\,\dd\vartheta\)、\(\theta^2=a\sin\vartheta\,\dd\varphi\)。由无挠结构方程得
\[
\omega^1{}_2=-\cos\vartheta\,\dd\varphi,\qquad
\Omega^1{}_2=\dd\omega^1{}_2
=\sin\vartheta\,\dd\vartheta\wedge\dd\varphi
=\frac1{a^2}\theta^1\wedge\theta^2.
\]
曲率二形式的系数即 \(K\)，故 \(K=a^{-2}\)。
\end{workedexample}
\par
\section{换用活动标架时怎样记录同一个曲率}

在非坐标标架 \(e_a\) 中，定义联络一形式
\(\omega^a{}_b(X)\) 使
\(\nabla_Xe_b=\omega^a{}_b(X)e_a\)。
把曲率定义作用于 \(e_b\)，使用一形式外微分的公式
\[
\dd\omega(X,Y)
=X[\omega(Y)]-Y[\omega(X)]-\omega([X,Y]),
\]
可整理为
\[
R(X,Y)e_b=\Omega^a{}_b(X,Y)e_a,\qquad
\Omega^a{}_b=\dd\omega^a{}_b+\omega^a{}_c\wedge\omega^c{}_b.
\]
第一项来自对联络系数求导，第二项来自连续两次基变换。
它与\cref{eq:math-curvature-components} 是同一张量的两种记法。

若 \(\theta^a\) 与 \(e_a\) 对偶，挠率形式
\(T^a(X,Y)=\theta^a(T(X,Y))\) 同样满足
\[
T^a=\dd\theta^a+\omega^a{}_b\wedge\theta^b.
\]
证明只需将右边作用于 \(e_b,e_c\)：\(\dd\theta^a\)
给出 \(-\theta^a([e_b,e_c])\)，另两项分别给出
\(\theta^a(\nabla_{e_b}e_c)\) 和
\(-\theta^a(\nabla_{e_c}e_b)\)，恰为挠率定义。
这说明标架不对易与几何挠率不是一回事：
Levi--Civita 联络虽无挠，非坐标标架仍可有 \(\dd\theta^a\ne0\)。

最后预告一个记法：把四指标 \(R^\ell{}_{kij}\) 缩并
\(R^k{}_{ikj}\) 得 Ricci 张量，再用 \(g^{ij}\) 缩并得标量曲率；
本书在后面讨论曲率变分时才用到，此处先不展开。




\section{标架与联络的公式联系}
结构函数描述标架的非完整性，联络形式描述平行移动；两者通过 Cartan 方程联系。下面汇集本章公式，并明确换标架时的变换方向。
\begin{note}[公式之间的联系]
\begin{align*}
[e_a,e_b]&=C^c{}_{ab}e_c,\\
\dd\theta^a&=-\frac12C^a{}_{bc}\theta^b\wedge\theta^c,\\
\nabla_{e_b}e_c&=\Gamma^a{}_{bc}e_a,\\
\omega^a{}_c&=\Gamma^a{}_{bc}\theta^b,\\
\Theta^a&=\dd\theta^a+\omega^a{}_b\wedge\theta^b,\\
\Omega^a{}_b&=\dd\omega^a{}_b+\omega^a{}_c\wedge\omega^c{}_b,\\
D\Theta^a&=\Omega^a{}_b\wedge\theta^b,\\
D\Omega^a{}_b&=0.
\end{align*}
若换标架写为 \(e' = e\Lambda\)，则
\[
\omega'=\Lambda^{-1}\omega\Lambda+\Lambda^{-1}\dd\Lambda,
\qquad
\Omega'=\Lambda^{-1}\Omega\Lambda.
\]
对应正文见\eqnrefs{eq:math-dg-frame-structure-equation}、\eqnrefs{eq:math-dg-connection-gauge-transform}、\eqnrefs{eq:math-dg-synthesis-cartan-first}、\eqnrefs{eq:math-dg-synthesis-cartan-second}与\eqnrefs{eq:math-dg-synthesis-bianchi}。
\end{note}

\begin{proposition}[局部联络的规范变换]
设主丛局部截面换为 $s^\prime=sg$，联络形式的局部代表为 $A=s^*\boldsymbol\omega$。则 $A^\prime=\operatorname{Ad}_{g^{-1}}A+g^{-1}\dd g$；矩阵群情形第一项为 $g^{-1}Ag$。
\end{proposition}
\begin{proof}
换局部截面 $s'=sg$ 时，新截面沿 $X$ 的切向量 $s'_*X$ 既含旧截面切向量的右平移，又含 $g$ 自身变化产生的垂直分量。联络形式对水平部分用 $G$-等变性、对垂直部分用再现性质，两项合起来给出非齐次变换律。

令 $s'=sg$，其中 $g:U\to G$。对 $X\in T_xM$，先计算新截面沿 $X$ 的切向量 $s'_*X$。由乘积截面微分，$s'_*X$ 同时含有 $s_*X$ 经右平移到 $s(x)g(x)$ 的水平部分，以及 $g(x)$ 自身沿 $X$ 变化产生的垂直部分。把后者经右平移回到单位元再经 基本 竖直 场 嵌入，得
\begin{equation*}
s'_*X
=(R_g)_*s_*X+\bigl(g^{-1}\dd g(X)\bigr)_P(sg),
\end{equation*}
其中 $(R_g)_*$ 是右平移 $R_g:h\mapsto hg$ 的切映射，$\bigl(g^{-1}\dd g(X)\bigr)_P$ 是把 $\mathfrak g$ 元素 $g^{-1}\dd g(X)$ 经 基本 场 映射 $\mathfrak g\to T_{sg}P$ 嵌入所得垂直向量。

由 $A'(X)=\boldsymbol\omega_{s'(x)}(s'_*X)$ 且 $s'(x)=s(x)g(x)$，代入上式并利用联络形式的线性性，
\begin{equation*}
A'(X)
=\boldsymbol\omega_{sg}\bigl((R_g)_*s_*X\bigr)
 +\boldsymbol\omega_{sg}\bigl(\bigl(g^{-1}\dd g(X)\bigr)_P(sg)\bigr).
\end{equation*}

对第一项用 $G$-等变性 $\boldsymbol\omega_{sg}\bigl((R_g)_*v\bigr)=\operatorname{Ad}_{g^{-1}}\boldsymbol\omega_s(v)$，取 $v=s_*X$，
\begin{equation*}
\boldsymbol\omega_{sg}\bigl((R_g)_*s_*X\bigr)
=\operatorname{Ad}_{g^{-1}}\boldsymbol\omega_s(s_*X)
=\operatorname{Ad}_{g^{-1}}A(X).
\end{equation*}
$G$-等变性是联络形式的定义公理之一，几何意义是右作用把联络形式按伴随表示拉回。

对第二项用再现性质 $\boldsymbol\omega_p\bigl(\xi_P(p)\bigr)=\xi$，取 $\xi=g^{-1}\dd g(X)$、$p=sg$，
\begin{equation*}
\boldsymbol\omega_{sg}\bigl(\bigl(g^{-1}\dd g(X)\bigr)_P(sg)\bigr)
=g^{-1}\dd g(X).
\end{equation*}
再现性质是联络形式的另一条定义公理：在基本竖直场上联络形式是恒等映射。

合并两项。
\begin{equation*}
A'(X)
=\operatorname{Ad}_{g^{-1}}A(X)+g^{-1}\dd g(X).
\end{equation*}
对矩阵 Lie 群 $\operatorname{Ad}_{g^{-1}}A=g^{-1}Ag$，于是
\begin{equation*}
A'(X)=g^{-1}A(X)g+g^{-1}\dd g(X),
\end{equation*}
即式\eqnrefs{eq:principal-local-connection-transform}。非齐次项来自局部截面本身也在变化，所以把联络系数当作普通张量会漏掉这一项。

典型例子。
\begin{itemize}
\item $G=U(1)$：$\operatorname{Ad}_{g^{-1}}=\mathrm{id}$（Abel 群 伴随 平凡），$A'=A+\ii\,\dd\chi$（取 $g=\ee^{\ii\chi}$）。这正是电磁势的规范变换 $A_\mu\mapsto A_\mu+\partial_\mu\chi$。
\item $G=SU(N)$：$A'=g^{-1}Ag+g^{-1}\dd g$，非 Abel 规范变换。$g^{-1}Ag$ 是共轭转动，$g^{-1}\dd g$ 是纯规范补偿。在 Yang--Mills 理论中，这一变换律保证 $D\psi=\dd\psi+A\psi$ 与 $\psi$ 同步齐次变换 $\psi\mapsto g^{-1}\psi$。
\end{itemize}

非齐次项 $g^{-1}\dd g$ 是规范势在规范变换下的"补偿"，对应物理上规范势 $A_\mu$ 在波函数相位重定向下获得的导数项。它使规范势不能视为张量，但保证协变导数 $D\psi$ 是张量（齐次变换）。这一非齐次性是规范理论全部相互作用结构的根源。
\end{proof}

\begin{proposition}[联络的换标架律]
对向量丛局部标架 $e^\prime=e\Lambda$，若 $\nabla e_b=e_a\otimes\omega^a{}_b$，则 $\omega^\prime=\Lambda^{-1}\omega\Lambda+\Lambda^{-1}\dd\Lambda$。
\end{proposition}
\begin{proof}
联络由 $\nabla e_c=e_a\otimes\omega^a{}_c$ 定义。换标架 $e'=e\Lambda$ 后，对 $\nabla e'_b$ 用 Leibniz 展开再换回新基，读出新系数即得非齐次变换律。

设新标架 $e'_b=e_c\Lambda^c{}_b$，其中 $\Lambda:U\to GL(n,\RR)$ 光滑。由联络对被求导向量满足 Leibniz 规则，对 $\nabla e'_b=\nabla(e_c\Lambda^c{}_b)$ 逐项展开，
\begin{align*}
\nabla e'_b
&=\nabla(e_c\Lambda^c{}_b)\\
&=(\nabla e_c)\Lambda^c{}_b+e_c\otimes\dd\Lambda^c{}_b.
\end{align*}
第一项是"基向量被联络作用"再乘系数，第二项是"系数被外微分"乘基向量——这正是联络作为方向导数的 Leibniz 性质。

由旧标架下联络一形式的定义 $\nabla e_c=e_a\otimes\omega^a{}_c$，代回第一项，
\begin{equation*}
(\nabla e_c)\Lambda^c{}_b
=e_a\otimes\omega^a{}_c\Lambda^c{}_b.
\end{equation*}

把指标 $c$ 在第一项中作哑指标，$a$ 作自由指标，合并两项
\begin{align*}
\nabla e'_b
&=e_a\otimes\omega^a{}_c\Lambda^c{}_b
 +e_c\otimes\dd\Lambda^c{}_b\\
&=e_a\otimes\bigl(\omega^a{}_c\Lambda^c{}_b+\dd\Lambda^a{}_b\bigr),
\end{align*}
其中第二行把第二项的哑指标 $c$ 改写为 $a$。

为读出新标架下的联络一形式 $\omega'$，需把基向量 $e_a$ 换回 $e'_r$。由 $e'_b=e_c\Lambda^c{}_b$ 取逆得 $e_a=e'_r(\Lambda^{-1})^r{}_a$，代回
\begin{align*}
\nabla e'_b
&=e'_r(\Lambda^{-1})^r{}_a\otimes
 \bigl(\omega^a{}_c\Lambda^c{}_b+\dd\Lambda^a{}_b\bigr)\\
&=e'_r\otimes
 \Bigl[
  (\Lambda^{-1})^r{}_a\omega^a{}_c\Lambda^c{}_b
  +(\Lambda^{-1})^r{}_a\dd\Lambda^a{}_b
 \Bigr].
\end{align*}

与新标架下定义 $\nabla e'_b=e'_r\otimes\omega'^r{}_b$ 比较基向量 $e'_r$ 的系数，得
\begin{equation*}
\omega'^r{}_b
=(\Lambda^{-1})^r{}_a\omega^a{}_c\Lambda^c{}_b
+(\Lambda^{-1})^r{}_a\dd\Lambda^a{}_b,
\end{equation*}
即\eqnrefs{eq:math-dg-connection-gauge-transform}。第一项是齐次的共轭转动，第二项是非齐次的纯规范补偿。

典型例子。
\begin{itemize}
\item 正交换标架 $\Lambda\in SO(n)$：$\Lambda^{-1}=\Lambda^T$，$\omega'=\Lambda^T\omega\Lambda+\Lambda^T\dd\Lambda$。第二项取值于 $\mathfrak{so}(n)$（反对称），保持度量相容性。
\item 局部 Lorentz 标架（四脚标架）变换 $\Lambda\in SO(1,3)$：$\omega'=\Lambda^{-1}\omega\Lambda+\Lambda^{-1}\dd\Lambda$，正是广义相对论中 自旋联络 在局部 Lorentz 变换下的规律。
\item 复标架变换 $\Lambda\in U(N)$：$\omega'=\Lambda^\dagger\omega\Lambda+\Lambda^\dagger\dd\Lambda$，对应 $U(N)$ 规范势的规范变换。
\end{itemize}

非齐次项 $\Lambda^{-1}\dd\Lambda$ 是"联络不是张量"的精确含义。在广义相对论中，它保证协变导数 $\nabla_\mu V^\nu=\partial_\mu V^\nu+\Gamma^\nu{}_{\mu\rho}V^\rho$ 在坐标变换下张量性成立——$\Gamma$ 的非齐次变换正好抵消 $\partial_\mu V^\nu$ 中多出的二阶项。在规范理论中，它使规范势 $A_\mu$ 在规范变换下获得 $\partial_\mu\chi$ 补偿项，从而协变导数 $D_\mu\psi$ 齐次变换。
\end{proof}

\begin{lemma}[协变外微分的平方]
设 $D$ 为向量丛联络诱导的协变外微分，曲率为 $\Omega$。对任意丛值微分形式 $\eta$，有 $D^2\eta=\Omega\wedge\eta$，右边使用曲率的纤维上的线性作用。
\end{lemma}
\begin{proof}
由协变外微分的定义 $D\eta=\dd\eta+\omega\wedge\eta$，其中 $\omega$ 是取值于 $\End E$ 的联络一形式、$\eta$ 是 $E$-值 $p$-形式。再次作用 $D$，先把 $D\eta$ 视作新的 $E$-值 $(p+1)$-形式代入同一公式，
\begin{equation*}
D^2\eta
=D(D\eta)
=\dd(D\eta)+\omega\wedge(D\eta).
\end{equation*}
把 $D\eta=\dd\eta+\omega\wedge\eta$ 代入两处，
\begin{equation*}
D^2\eta
=\dd(\dd\eta+\omega\wedge\eta)
 +\omega\wedge(\dd\eta+\omega\wedge\eta).
\end{equation*}
由外微分的线性性拆开，
\begin{equation*}
D^2\eta
=\dd^2\eta+\dd(\omega\wedge\eta)
 +\omega\wedge\dd\eta+\omega\wedge(\omega\wedge\eta).
\end{equation*}
由 $\dd^2=0$ 第一项消失。对第二项用分次 Leibniz 规则 $\dd(\alpha\wedge\beta)=\dd\alpha\wedge\beta+(-1)^{|\alpha|}\alpha\wedge\dd\beta$，其中 $|\omega|=1$，
\begin{equation*}
\dd(\omega\wedge\eta)
=\dd\omega\wedge\eta+(-1)^1\omega\wedge\dd\eta
=\dd\omega\wedge\eta-\omega\wedge\dd\eta.
\end{equation*}
代回，
\begin{equation*}
D^2\eta
=\dd\omega\wedge\eta-\omega\wedge\dd\eta
 +\omega\wedge\dd\eta+\omega\wedge\omega\wedge\eta.
\end{equation*}
中间两项 $-\omega\wedge\dd\eta+\omega\wedge\dd\eta$ 严格相消，只剩
\begin{equation*}
D^2\eta
=\dd\omega\wedge\eta+\omega\wedge\omega\wedge\eta
=(\dd\omega+\omega\wedge\omega)\wedge\eta.
\end{equation*}
由曲率二形式定义 $\Omega:=\dd\omega+\omega\wedge\omega$，代回得
\begin{equation*}
D^2\eta=\Omega\wedge\eta,
\end{equation*}
即\eqnrefs{eq:math-dg-covariant-derivative-square}。这也解释了为何曲率必然以 $\dd\omega+\omega\wedge\omega$ 而不是单独 $\dd\omega$ 出现。
\end{proof}

\begin{proposition}[曲率的有限扰动]
若 $\widetilde\nabla=\nabla+A$ 且 $A\in\Omega^1(M;\End E)$，则 $\widetilde R=R+D_\nabla A+A\wedge A$。此式是有限差的精确恒等式，未截断扰动阶。
\end{proposition}
\begin{proof}
曲率是联络的二次函数（含 $\nabla^2$ 和 $\nabla\wedge\nabla$），故联络的有限扰动 $A$ 产生线性项（$D_\nabla A$）和二次项（$A\wedge A$），结构完全由交换子代数决定。

把新联络写成
\begin{equation*}
 \widetilde\nabla_X=\nabla_X+A_X,
 \qquad
 A_X:=A(X,\cdot).
\end{equation*}
$A$ 是 $\Omega^1(M;\operatorname{End}E)$ 中任一元素，$\widetilde\nabla$ 与 $\nabla$ 同为 $E$ 上联络当且仅当它们的差是张量一形式。

曲率是交换子减联络沿括号方向的作用：
\begin{equation*}
 \widetilde R(X,Y)
 =[\widetilde\nabla_X,\widetilde\nabla_Y]-\widetilde\nabla_{[X,Y]}
 =[\nabla_X+A_X,\nabla_Y+A_Y]
 -(\nabla_{[X,Y]}+A_{[X,Y]}).
\end{equation*}

按双线性展开：
\begin{itemize}
\item 纯 $\nabla$ 部分：$[\nabla_X,\nabla_Y]-\nabla_{[X,Y]}=R(X,Y)$，即旧曲率。
\item 一次 $A$ 的交叉项：$\nabla_XA_Y-\nabla_YA_X-A_{[X,Y]}$，组合成协变外微分 $(D_\nabla A)(X,Y)$，即 $A$ 沿旧联络的协变外微分。
\item 二次项：$[A_X,A_Y]=A_XA_Y-A_YA_X$，即 $A\wedge A$。
\end{itemize}
逐项核对：$[\nabla_X,A_Y]=\nabla_XA_Y-A_Y\nabla_X$ 中 $-A_Y\nabla_X$ 与 $[A_X,\nabla_Y]=A_X\nabla_Y-\nabla_YA_X$ 中的 $A_X\nabla_Y$ 不直接相消，但与 $\nabla_XA_Y-\nabla_YA_X$ 合并后正好是协变外微分 $D_\nabla A$ 的定义。

因此
\begin{equation*}
 \widetilde R=R+D_\nabla A+A\wedge A,
\end{equation*}
与\eqnrefs{eq:math-dg-curvature-finite-perturbation}完全相同。曲率对联络"线性化 + 自相互作用"的结构因此不是坐标巧合，而是联络仿射几何的必然结果：线性项 $D_\nabla A$ 是一阶变分（用于 Chern--Weil 理论和指标定理），二次项 $A\wedge A$ 是非 Abel 自相互作用（与 Yang--Mills 方程的非线性性同源）。

取 $A_t=tA$，$\widetilde R_t=R+tD_\nabla A+t^2 A\wedge A$。对 $\operatorname{tr}(\widetilde R_t^k)$ 在 $t=0$ 求导，只有线性项 $D_\nabla A$ 贡献（二次项 $t^2$ 在 $t=0$ 导数为零）。这正是一阶变分控制 Chern--Weil 超渡 的原因：特征类的变分只看 $D_\nabla A$，非 Abel 自相互作用 $A\wedge A$ 在一阶变分中不可见，但在二阶变分与 Yang--Mills 作用量中起决定作用。

Yang--Mills 作用量 $S=\int\operatorname{tr}(F\wedge *F)$ 在 $\widetilde\nabla=\nabla+A$ 下展开为
\begin{equation*}
 S(\widetilde\nabla)=S(\nabla)+2\int\operatorname{tr}(D_\nabla A\wedge *F)+2\int\operatorname{tr}(A\wedge A\wedge *F)+\int\operatorname{tr}((D_\nabla A+A\wedge A)\wedge *(D_\nabla A+A\wedge A)).
\end{equation*}
第一阶变分给 Yang--Mills 方程 $D_\nabla *F=0$；二阶变分含 $A\wedge A$ 的二次与三次项，是规范场自能的非线性来源。Abel 情形 $A\wedge A=0$，作用量退化为 Maxwell 二次型，变分方程线性。

典型例子。
\begin{itemize}
\item $U(1)$ 规范理论（Abel）：$A\wedge A=0$（交换子为零），$\widetilde F=F+\dd A$。电磁场强度在规范势平移下只加 $\dd A$，没有自相互作用——这正是 Maxwell 方程线性的根源。
\item $SU(2)$ Yang--Mills：$\widetilde F=F+D_\nabla A+A\wedge A$，$A\wedge A\neq0$ 产生非 Abel 自相互作用。Yang--Mills 方程 $D_\mu F^{\mu\nu}=0$ 展开后含 $A_\mu A_\nu A^\nu$ 三次项，是非线性性的来源。
\item Riemann 几何：取 $\nabla$ 为 Levi--Civita 联络、$A$ 为 $(1,2)$ 张量，$\widetilde R=R+\nabla A+A\circ A$ 给出两个度量相容联络的曲率差。若 $\widetilde\nabla$ 也无挠，$A$ 对下指标对称，二次项控制如 Rauch 比较定理中的曲率比较。
\end{itemize}

这一展开是规范理论全部相互作用结构的微缩版：线性项 $D_\nabla A$ 是"外源激发"，二次项 $A\wedge A$ 是"自相互作用"。在 Yang--Mills 中，规范场既被外源激发又自我激发，正是非 Abel 规范理论区别于 Maxwell 理论的本质。在广义相对论中，Einstein 张量 $G_{\mu\nu}$ 含 Ricci 曲率，而 Ricci 在联络扰动下的二次项对应引力场的自相互作用——这是 Einstein 方程非线性的几何根源。
\end{proof}

\begin{theorem}[Cartan 结构方程]
设 $\theta^a$ 为切丛余标架，$\omega^a{}_b$ 为联络形式。挠率与曲率二形式分别满足 $T^a=\dd\theta^a+\omega^a{}_b\wedge\theta^b$、$\Omega^a{}_b=\dd\omega^a{}_b+\omega^a{}_c\wedge\omega^c{}_b$。
\end{theorem}
\begin{proof}
先从挠率定义出发。对任意 $X,Y$，
\[
\begin{aligned}
T^a(X,Y)
={}&\theta^a(\nabla_XY-\nabla_YX-[X,Y])\\
={}&X[\theta^a(Y)]-Y[\theta^a(X)]-\theta^a([X,Y])\\
&+\omega^a{}_b(X)\theta^b(Y)
-\omega^a{}_b(Y)\theta^b(X).
\end{aligned}
\]
第一行中前三项恰是 $\dd\theta^a(X,Y)$，后二项恰是
$(\omega^a{}_b\wedge\theta^b)(X,Y)$。故
\begin{equation}\label{eq:math-dg-synthesis-cartan-first}
{T^a=\dd\theta^a+\omega^a{}_b\wedge\theta^b.}
\end{equation}

再对基向量 $e_b$ 展开曲率定义：
\[
\begin{aligned}
R(X,Y)e_b
={}&\nabla_X\bigl(\omega^c{}_b(Y)e_c\bigr)
-\nabla_Y\bigl(\omega^c{}_b(X)e_c\bigr)
-\omega^c{}_b([X,Y])e_c\\
={}&\Bigl(
X[\omega^a{}_b(Y)]-Y[\omega^a{}_b(X)]-\omega^a{}_b([X,Y])
\Bigr)e_a\\
&+\Bigl(
\omega^c{}_b(Y)\omega^a{}_c(X)
-\omega^c{}_b(X)\omega^a{}_c(Y)
\Bigr)e_a.
\end{aligned}
\]
把第一括号认作 $\dd\omega^a{}_b(X,Y)$，第二括号认作
$(\omega^a{}_c\wedge\omega^c{}_b)(X,Y)$，得到
\begin{equation}\label{eq:math-dg-synthesis-cartan-second}
{
\Omega^a{}_b
=\dd\omega^a{}_b+\omega^a{}_c\wedge\omega^c{}_b.}
\end{equation}

第一结构方程 $T^a=\dd\theta^a+\omega^a{}_b\wedge\theta^b$ 说挠率是「外微分加联络修正」；第二结构方程 $\Omega^a{}_b=\dd\omega^a{}_b+\omega^a{}_c\wedge\omega^c{}_b$ 说曲率是联络的「外微分加自乘」。二者合起来把挠率和曲率完全用标架 $\{\theta^a\}$ 和联络 $\{\omega^a{}_b\}$ 表达，无需坐标。对 Riemann 流形（挠率 $T=0$），第一方程显式解出 $\omega^a{}_b$，将其表为 $\dd\theta^a$ 的函数，第二方程则给出曲率——这是 Cartan 活动标架法计算曲率的标准流程。

取球面 $\mathbb S^2$ 的活动标架 $\theta^1=R\,\dd\theta$、$\theta^2=R\sin\theta\,\dd\varphi$。由 $T=0$ 解第一方程得 $\omega^1{}_2=-\cos\theta\,\dd\varphi$。代入第二方程：$\Omega^1{}_2=\dd\omega^1{}_2=\sin\theta\,\dd\theta\wedge\dd\varphi=(1/R^2)\,\theta^1\wedge\theta^2$，立即读出截面曲率 $K=1/R^2$。整个计算无需 Christoffel 符号。
\end{proof}

\begin{proposition}[曲率的齐次规范变换]
若 $A^\prime=g^{-1}Ag+g^{-1}\dd g$，则 $F^\prime=g^{-1}Fg$。因此任何伴随不变多项式作用于曲率所得的标量值形式均不依赖局部标架。
\end{proposition}
\begin{proof}
规范变换下联络形式的变换律为
\begin{equation*}
 A'=g^{-1}Ag+g^{-1}\dd g.
\end{equation*}
记 $\eta=g^{-1}\dd g$（纯规范），则 $A'=g^{-1}Ag+\eta$。目标是对 $F'=\dd A'+A'\wedge A'$ 直接计算并证明 $F'=g^{-1}Fg$。

由 $\eta=g^{-1}\dd g$ 是纯规范，Maurer--Cartan 方程给出
\begin{equation*}
 \dd\eta+\eta\wedge\eta=0.
\end{equation*}
这是 $\eta$ 的"平坦性"条件，几何意义是 $g$ 是全局定义的光滑函数，其微分可积。

由 $\dd(g^{-1})=-\eta\,g^{-1}$（对 $g\cdot g^{-1}=I$ 求微分），
\begin{equation*}
 \dd(g^{-1}Ag)
 =g^{-1}(\dd A)g
 -\eta\wedge g^{-1}Ag
 -g^{-1}Ag\wedge\eta.
\end{equation*}
后两项是 $\eta$ 与 $g^{-1}Ag$ 的交叉楔积，来自 $\dd(g^{-1})$ 与 $\dd g$ 上的 Leibniz 规则。

展开 $F'=\dd A'+A'\wedge A'$ 后收集所有项：
\begin{itemize}
\item $g^{-1}(\dd A)g$：来自 $\dd(g^{-1}Ag)$。
\item $g^{-1}(A\wedge A)g$：来自 $(g^{-1}Ag)\wedge(g^{-1}Ag)$。
\item $\dd\eta$：来自 $\dd\eta$。
\item $\eta\wedge\eta$：来自 $\eta\wedge\eta$。
\item 交叉项 $\eta\wedge g^{-1}Ag$ 和 $g^{-1}Ag\wedge\eta$：来自 $\dd(g^{-1}Ag)$ 展开的后两项和 $A'\wedge A'$ 的展开。
\end{itemize}
具体地，$A'\wedge A'=(g^{-1}Ag+\eta)\wedge(g^{-1}Ag+\eta)$ 展开为四项：$(g^{-1}Ag)\wedge(g^{-1}Ag)$、$(g^{-1}Ag)\wedge\eta$、$\eta\wedge(g^{-1}Ag)$、$\eta\wedge\eta$。其中两个交叉项与 $\dd(g^{-1}Ag)$ 中的两个交叉项符号相反，精确配对抵消。

由 Maurer--Cartan 恒等式 $\dd\eta=-\eta\wedge\eta$，以及交叉项精确配对抵消，只剩
\begin{equation*}
 F'=g^{-1}(\dd A+A\wedge A)g
 =g^{-1}Fg.
\end{equation*}
交叉项的相消可显式核对：$\dd(g^{-1}Ag)$ 给 $-\eta\wedge g^{-1}Ag-g^{-1}Ag\wedge\eta$，$A'\wedge A'$ 给 $+g^{-1}Ag\wedge\eta+\eta\wedge g^{-1}Ag$，二者符号相反；$\dd\eta$ 与 $\eta\wedge\eta$ 则由 Maurer--Cartan 恒等式消去。

典型例子。
\begin{itemize}
\item $U(1)$ 规范：$g=\ee^{\ii\chi}$，$F'=F$（Abel 群共轭平凡）。电磁场强度 $F_{\mu\nu}=\partial_\mu A_\nu-\partial_\nu A_\mu$ 在规范变换 $A_\mu\mapsto A_\mu+\partial_\mu\chi$ 下不变，正是 Maxwell 理论规范不变性的几何根源。
\item $SU(N)$ Yang--Mills：$F'=g^{-1}Fg$，规范场强做共轭变换。$\operatorname{tr}(F\wedge F)$ 因此规范不变，是 Chern--Weil 理论构造第二 Chern 类 $c_2(E)$ 的基础。
\item $SO(n)$ 标架丛：$\Omega'=g^{-1}\Omega g$ 给 Riemann 曲率二形式的齐次变换律，与联络一形式的非齐次律对比鲜明。
\item 仿射联络：$\Gamma$ 在坐标变换下非齐次（含二阶导数项），但 Riemann 曲率张量 $R^\rho{}_{\sigma\mu\nu}$ 齐次变换，是真正的张量。
\item 瞬子解：$SU(2)$ BPST 瞬子 $A_\mu=(1-f)\,g^{-1}\partial_\mu g$ 中 $F=g^{-1}F_0g$ 共轭不变，$\operatorname{tr}(F\wedge F)=\operatorname{tr}(F_0\wedge F_0)$ 给瞬子数 $k=1$。
\end{itemize}

这就是"联络 非张量而 曲率 张量"的矩阵形式证明：$A$ 在 规范 变换下有非齐次项 $g^{-1}\dd g$（仿射变换），但 $F$ 只做共轭变换（齐次变换），因此曲率是真正的张量。在 Yang--Mills 理论中，作用量 $\int\operatorname{tr}(F\wedge *F)$ 规范不变正是源于 $F$ 的齐次变换律与 迹 的循环性。在 Chern--Weil 理论中，$\operatorname{tr}(F^k)$ 规范不变且闭，定义了不依赖联络选择的拓扑不变量。这一"非齐次 $\to$ 齐次"的升级是规范理论全部构造性结论（守恒流、规范不变作用量、拓扑荷）的共同根源。
\end{proof}

\begin{theorem}[Bianchi 恒等式]
由同一切丛联络定义的挠率 $T$ 与曲率 $\Omega$ 满足 $DT^a=\Omega^a{}_b\wedge\theta^b$ 和 $D\Omega=0$；第二式也适用于一般向量丛联络。
\end{theorem}
\begin{proof}
由第一结构方程，
\[
T^a=\dd\theta^a+\omega^a{}_b\wedge\theta^b.
\]
用\eqnrefs{eq:math-dg-exterior-covariant-derivative}计算
\[
\begin{aligned}
D T^a
={}&\dd T^a+\omega^a{}_c\wedge T^c\\
={}&\dd\omega^a{}_b\wedge\theta^b
-\omega^a{}_b\wedge\dd\theta^b
+\omega^a{}_c\wedge\dd\theta^c
+\omega^a{}_c\wedge\omega^c{}_b\wedge\theta^b.
\end{aligned}
\]
中间两项在重命名哑指标后相消，于是
\[
D T^a
=\bigl(\dd\omega^a{}_b+\omega^a{}_c\wedge\omega^c{}_b\bigr)\wedge\theta^b
=\Omega^a{}_b\wedge\theta^b.
\]
因此
\begin{equation}\label{eq:math-dg-bianchi-first}
{D T^a=\Omega^a{}_b\wedge\theta^b.}
\end{equation}

对第二结构方程作协变外微分。对一个 $\End(TM)$-值二形式，
\[
(D\Omega)^a{}_b
=\dd\Omega^a{}_b
+\omega^a{}_c\wedge\Omega^c{}_b
-\Omega^a{}_c\wedge\omega^c{}_b.
\]
代入 $\Omega=\dd\omega+\omega\wedge\omega$，并使用 $\dd^2=0$，
\[
\begin{aligned}
D\Omega
={}&\dd(\omega\wedge\omega)
+\omega\wedge(\dd\omega+\omega\wedge\omega)
-(\dd\omega+\omega\wedge\omega)\wedge\omega\\
={}&(\dd\omega\wedge\omega-\omega\wedge\dd\omega)
+\omega\wedge\dd\omega+\omega\wedge\omega\wedge\omega\\
&-\dd\omega\wedge\omega-\omega\wedge\omega\wedge\omega\\
={}&0.
\end{aligned}
\]
故
\begin{equation}\label{eq:math-dg-bianchi-second}
{D\Omega^a{}_b=0.}
\end{equation}
两式也可合写为
\begin{equation}\label{eq:math-dg-synthesis-bianchi}
D T^a=\Omega^a{}_b\wedge\theta^b,
\qquad
D\Omega^a{}_b=0.
\end{equation}

Bianchi 恒等式是曲率的可积条件：$D\Omega=0$ 说明曲率二形式在外协变微分下闭合。这是 Chern--Weil 理论的出发点：$\operatorname{tr}(\Omega^k)$ 为闭形式，其联络无关性由下一定理证明。二维曲面上的 $D\Omega$ 是三形式，因维数自动为零，并不要求 Gauss 曲率为常数。规范理论的 Bianchi 恒等式同样由联络的曲率定义得到；Einstein 张量的 $\nabla_\mu G^{\mu\nu}=0$ 则来自 Levi--Civita 联络第二 Bianchi 恒等式的指标缩并，不能由独立 Yang--Mills 联络的恒等式推出。
\end{proof}

\begin{theorem}[Chern--Weil 形式的联络无关性]
固定一个复向量丛，令 $A_t$ 为光滑联络族、$F_t$ 为曲率。对 $k\ge1$，$\operatorname{tr}(F_t^k)$ 闭合，且 $\partial_t\operatorname{tr}(F_t^k)=k\dd\operatorname{tr}(\dot A_t\wedge F_t^{k-1})$。因此其 de Rham 类不随联络变化。
\end{theorem}
\begin{proof}
对一族联络 $A_t$ 与曲率 $F_t$，把 $\operatorname{tr}(F_t^k)$ 对 $t$ 求导。利用 $\dot F_t=D_t\dot A_t$、Bianchi 恒等式 $D_tF_t=0$ 与 迹 的循环性，把导数写成 恰当形式，从而 de Rham 上同调类与 $t$ 无关。

对
\begin{equation*}
F_t=\dd A_t+A_t\wedge A_t
\end{equation*}
求导（$\dot A_t:=\frac{\dd A_t}{\dd t}$）：
\begin{equation*}
\dot F_t
=\dd\dot A_t+A_t\wedge\dot A_t+
\dot A_t\wedge A_t
=D_t\dot A_t.
\end{equation*}
最后等式是协变外微分的定义 $D_t\dot A_t=\dd\dot A_t+[A_t,\dot A_t]$，其中 $[A_t,\dot A_t]=A_t\wedge\dot A_t+\dot A_t\wedge A_t$（一形式楔积带符号合并）。

由 Leibniz 规则，
\begin{equation*}
\frac{\dd}{\dd t}\operatorname{tr}(F_t^k)
=k\,\operatorname{tr}(\dot F_t\wedge F_t^{k-1})
=k\,\operatorname{tr}(D_t\dot A_t\wedge F_t^{k-1}).
\end{equation*}
因子 $k$ 来自 $F_t^k=F_t\wedge\cdots\wedge F_t$ 中 $\dot F_t$ 可落在 $k$ 个位置之一，迹 的循环性使这 $k$ 项相等。

利用协变 Leibniz 规则和 Bianchi 恒等式，可把 $D_t$ 从 $\dot A_t$ 移到 $F_t^{k-1}$：
\begin{equation*}
D_t(\dot A_t\wedge F_t^{k-1})
=D_t\dot A_t\wedge F_t^{k-1}
-\dot A_t\wedge D_t(F_t^{k-1}),
\end{equation*}
以及第二 Bianchi 恒等式 $D_tF_t=0$（从而 $D_t(F_t^{k-1})=0$），得
\begin{equation*}
D_t(\dot A_t\wedge F_t^{k-1})
=D_t\dot A_t\wedge F_t^{k-1}.
\end{equation*}
这一步是关键：Bianchi 恒等式让 $D_t$ 在乘积上"穿过"所有 $F_t$ 因子，仿佛它们是"协变常量"。

迹把 $D_t$ 退化为 $\dd$。对任意次数为 $p$ 的矩阵值形式 $\beta$，分次循环性给出 $\operatorname{tr}(A_t\wedge\beta)=(-1)^p\operatorname{tr}(\beta\wedge A_t)$，故分次交换子 $A_t\wedge\beta-(-1)^p\beta\wedge A_t$ 的迹为零。因此
\begin{equation*}
\operatorname{tr}(D_t(\dot A_t\wedge F_t^{k-1}))
=\dd\operatorname{tr}(\dot A_t\wedge F_t^{k-1}),
\end{equation*}
得到\eqnrefs{eq:chern-weil-transgression}。

积分 $t\in[0,1]$ 后，
\begin{equation*}
\operatorname{tr}(F_1^k)-\operatorname{tr}(F_0^k)
=k\,\dd\int_0^1\operatorname{tr}(\dot A_t\wedge F_t^{k-1})\,\dd t,
\end{equation*}
两组 联络 产生的 特征 形式 相差 恰当形式，因此代表同一 de Rham 类 $[\operatorname{tr}(F^k)]\in H_{\mathrm{dR}}^{2k}(M)$。

典型例子。
\begin{itemize}
\item $U(1)$ 线丛：$k=1$，$\operatorname{tr}(F)=F$ 是闭二形式，$[\ii F/(2\pi)]\in H^2_{\mathrm{dR}}(M)$ 是第一 Chern 类 $c_1(E)$。两个不同联络的 $F$ 相差 恰当形式，但积分 $\frac{\ii}{2\pi}\int_M F$ 给出 Chern 数（如磁通量子化）。
\item $SU(N)$ 丛：$k=2$，$\operatorname{tr}(F\wedge F)$ 是闭四形式，$[\operatorname{tr}(F\wedge F)/8\pi^2]$ 是第二 Chern 类 $c_2(E)$。Yang--Mills 瞬子数 $k=\int\operatorname{tr}(F\wedge F)/8\pi^2$ 是拓扑不变量，不依赖规范选择。
\end{itemize}

Chern--Weil 理论给出从 $\mathfrak g$ 上的不变多项式到 $H_{\mathrm{dR}}^*(M)$ 的同态：每个 $G$-不变多项式 $P$ 给出闭形式 $P(F)$，其 de Rham 类只依赖丛 $E$ 的拓扑。结合 Atiyah--Singer 指标定理，这些特征类积分给出椭圆算子的解析指标——如 Dirac 算子的 $\hat A$ 类、符号丛的 Chern 字符。这是"曲率局部数据决定全局拓扑"的精确机制。
\end{proof}

\begin{proposition}[路径序指数的收敛构造]
设 $A_t$ 为有限闭区间上的连续矩阵函数。初值问题 $\dot V=-A_tV$ 的解为绝对一致收敛的 Dyson 级数。若所有 $A_t$ 两两对易，则解退化为 $V(t)=\exp(-\int_0^tA_s\,\dd s)V(0)$。
\end{proposition}
\begin{proof}
把\eqnrefs{eq:parallel-matrix-ode} 即 $\dot V(t)=-A_tV(t)$、$V(0)$ 给定，从 $0$ 到 $t$ 积分得积分方程
\begin{equation*}
V(t)-V(0)=-\int_0^tA_{t_1}V(t_1)\,\dd t_1,
\end{equation*}
即
\begin{equation*}
V(t)=V(0)-\int_0^tA_{t_1}V(t_1)\,\dd t_1.
\end{equation*}
把右端 $V(t_1)$ 用同一积分方程代入，$V(t_1)=V(0)-\int_0^{t_1}A_{t_2}V(t_2)\,\dd t_2$，得
\begin{align*}
V(t)
&=V(0)-\int_0^tA_{t_1}
 \Bigl[V(0)-\int_0^{t_1}A_{t_2}V(t_2)\,\dd t_2\Bigr]\dd t_1\\
&=V(0)-\int_0^tA_{t_1}\dd t_1\,V(0)
 +\int_0^t\int_0^{t_1}A_{t_1}A_{t_2}V(t_2)\,\dd t_2\dd t_1\\
&=\Bigl[I-\int_0^tA_{t_1}\dd t_1\Bigr]V(0)
 +\int_0^t\int_0^{t_1}A_{t_1}A_{t_2}V(t_2)\,\dd t_2\dd t_1.
\end{align*}
对余项中的 $V(t_2)$ 再代入一次，$V(t_2)=V(0)-\int_0^{t_2}A_{t_3}V(t_3)\,\dd t_3$，
\begin{align*}
\int_0^t\int_0^{t_1}A_{t_1}A_{t_2}V(t_2)\,\dd t_2\dd t_1
&=\int_0^t\int_0^{t_1}A_{t_1}A_{t_2}\dd t_2\dd t_1\,V(0)\\
&\quad-\int_0^t\int_0^{t_1}\int_0^{t_2}
 A_{t_1}A_{t_2}A_{t_3}V(t_3)\,\dd t_3\dd t_2\dd t_1.
\end{align*}
合并得第三次迭代
\begin{align*}
V(t)
&=\Bigl[I-\int_0^tA_{t_1}\dd t_1
 +\int_0^t\int_0^{t_1}A_{t_1}A_{t_2}\dd t_2\dd t_1\Bigr]V(0)\\
&\quad-\int_0^t\int_0^{t_1}\int_0^{t_2}
 A_{t_1}A_{t_2}A_{t_3}V(t_3)\,\dd t_3\dd t_2\dd t_1.
\end{align*}
由归纳法，截断至 $N$ 阶的有限展开为
\begin{equation*}
V(t)=\sum_{j=0}^{N}(-1)^j
\int_{0<t_j<\cdots<t_1<t}
A_{t_1}\cdots A_{t_j}\,\dd t_j\cdots\dd t_1\,V(0)+R_{N+1}(t),
\end{equation*}
其中 $j=0$ 项理解为 $I\,V(0)$，余项为
\[
R_{N+1}(t)=(-1)^{N+1}\int_{0<t_{N+1}<\cdots<t_1<t}
A_{t_1}\cdots A_{t_{N+1}}V(t_{N+1})\,\dd t_{N+1}\cdots\dd t_1.
\]
在有限区间 $0\le t\le T$ 上令 $M=\sup_t\|A_t\|$、$B=\sup_t\|V(t)\|<\infty$，则单纯形体积给出
\[
\|R_{N+1}(t)\|\le B\frac{M^{N+1}T^{N+1}}{(N+1)!}\longrightarrow0,
\qquad
\|\text{第 }j\text{ 阶项}\|\le\|V(0)\|\frac{M^jT^j}{j!}.
\]
这分别证明余项一致趋零及级数绝对一致收敛，从而产生 Dyson 级数
\begin{equation*}
V(t)=\Bigl[I+\sum_{k=1}^\infty(-1)^k
\int_{0<t_k<\cdots<t_1<t}
A_{t_1}\cdots A_{t_k}\,\dd t_k\cdots\dd t_1\Bigr]V(0).
\end{equation*}
若所有 $A_t$ 对易，积分区域的 $k!$ 个排列给相同 integrand，于是有序单纯形积分等于立方体积分的 $1/k!$，级数便退化为普通指数。这解释了什么时候可以安全地把 $\mathcal P$ 去掉。
\end{proof}

